% !TeX TXS-program:compile = txs:///arara
% arara: pdflatex: {shell: no, synctex: no, interaction: batchmode}
% arara: pdflatex: {shell: no, synctex: no, interaction: batchmode}

\documentclass[11pt,a4paper]{ltxdoc}
\usepackage{bera}
\usepackage{inconsolata}
\usepackage[T1]{fontenc}
\usepackage[scale=0.875]{cabin}
\usepackage[bm]{mathador}
\usepackage{fancyvrb}
\usepackage{fancyhdr}
\usepackage{tabularray}
\usepackage{fontawesome5}
\fancyhf{}
\renewcommand{\headrulewidth}{0pt}
\lfoot{\sffamily\small [mathador]}
\cfoot{\sffamily\small - \thepage{} -}
\rfoot{\hyperlink{matoc}{\small\faArrowAltCircleUp[regular]}}
\usepackage{hologo}
\providecommand\tikzlogo{Ti\textit{k}Z}
\providecommand\TeXLive{\TeX{}Live\xspace}
\let\TikZ\tikzlogo

\usepackage{hyperref}
\urlstyle{same}
\hypersetup{pdfborder=0 0 0}
\usepackage[margin=2cm]{geometry}
\setlength{\parindent}{0pt}
\def\TPversion{0.1.1}
\def\TPdate{25/09/2025}
\usepackage{tcolorbox}
\tcbuselibrary{breakable,skins,hooks,listingsutf8}
%\usepackage{soul}
%\sethlcolor{lightgray!25}

\lstset{
	language=[LaTeX]TeX,%
	basicstyle=\ttfamily,%
	keywordstyle={\color{blue}},%
	classoffset=0,%
	keywords={},%
	alsoletter={-},%
	keywordstyle={\color{blue}},%
	classoffset=1,%
	alsoletter={-},%
	morekeywords={mathador},%
	keywordstyle={\color{violet}},%
	classoffset=2,%
	alsoletter={-},%
	morekeywords={\PlateauMathador,\DeMathador,\CibleMathador,\TirageDesMathador},%
	keywordstyle={\color{green!50!black}},%
	classoffset=3,%
	morekeywords={},%
	keywordstyle={\color{orange}}
}

\newtcblisting{DemoCode}[1]{%
	enhanced,width=\linewidth,%
	bicolor,size=title,%
	colback=cyan!10!white,%
	colbacklower=cyan!5!white,%
	colframe=cyan!75!black,%
	listing options={%
		breaklines=true,%
		breakatwhitespace=true,%
		style=tcblatex,basicstyle=\small\ttfamily,%
		tabsize=4,%
		commentstyle={\itshape\color{gray}},
		keywordstyle={\color{blue}},%
		classoffset=0,%
		keywords={\tikzset,\usepackage,\definecolor,\includegraphics,xstring,listofitems,\readlist,\showitems,\xintFor,\xintSeq},%
		alsoletter={-},%
		keywordstyle={\color{blue}},%
		classoffset=1,%
		alsoletter={-},%
		morekeywords={mathador},%
		keywordstyle={\color{violet}},%
		classoffset=2,%
		alsoletter={-},%
		morekeywords={\PlateauMathador,\DeMathador,\CibleMathador,\TirageDesMathador},%
		keywordstyle={\color{green!50!black}},%
		classoffset=3,%
		morekeywords={},%
		keywordstyle={\color{orange}}
	},%
	#1
}

\newtcbinputlisting\DemoCodeFile[1]{%
	enhanced,width=\linewidth,%
	bicolor,size=title,%
	colback=lightgray!10!white,%
	colbacklower=lightgray!5!white,%
	colframe=lightgray!75!black,%
	listing options={%
		breaklines=true,%
		breakatwhitespace=true,%
		style=tcblatex,
		basicstyle=\scriptsize\ttfamily,%
		tabsize=4,%
		commentstyle={\itshape\color{gray}},%
		%lastline=148
	},%
	breakable,
	listing only,%
	listing file={#1}
}

\NewDocumentCommand\ShowCode{ m }{%
	\lstinline{#1}%
}

\usepackage[french]{babel}

\begin{document}

\thispagestyle{empty}

\begin{center}
	\begin{minipage}{0.88\linewidth}
		\begin{tcolorbox}[colframe=yellow,colback=yellow!15]
			\begin{center}
				\renewcommand{\arraystretch}{1.25}%
				\begin{tabular}{c}
					{\Huge \texttt{mathador}}\\
					\\
					{\LARGE De commandes pour travailler avec le} \\
					{\LARGE jeu Mathador : plateau / dés / tirage.} \\
					\\
					{\small \texttt{Version \TPversion{} -- \TPdate}}
				\end{tabular}
			\end{center}
		\end{tcolorbox}
	\end{minipage}
\end{center}

\begin{center}
	\begin{tabular}{c}
		\texttt{Cédric Pierquet} : {\ttfamily c pierquet -- at -- outlook . fr}\\
		\texttt{\url{https://github.com/cpierquet/latex-packages/tree/main//mathador}} \\
	\end{tabular}
\end{center}

\hrule

\begin{center}
	\begin{tabular}{c}
		{\ttfamily Auteur du jeu :  Éric Trouillot}\\
		{\ttfamily Éditeur du jeu :  Réseau Canopé (\url{https://www.mathador.fr})}
	\end{tabular}
\end{center}

\hrule

\vfill

\begin{tcblisting}{colframe=lightgray,colback=lightgray!5,listing only}
\PlateauMathador{1,6,19,7,12}{39}
\end{tcblisting}

\begin{tcolorbox}[colframe=lightgray,colback=lightgray!5]
\hfill\PlateauMathador{1,6,19,7,12}{39}\hfill\null
\end{tcolorbox}

\begin{tcblisting}{colframe=lightgray,colback=lightgray!5,listing only}
\TirageDesMathador{39}{1}{6}{7}{12}{19}
\end{tcblisting}

\begin{tcolorbox}[colframe=lightgray,colback=lightgray!5]
\hfill\TirageDesMathador{39}{1}{6}{7}{12}{19}\hfill\null
\end{tcolorbox}

\vfill\null

\pagebreak

\phantomsection

\hypertarget{matoc}{}

\tableofcontents

\vspace*{5mm}

%\hrule

\pagebreak

\section{Introduction}

\subsection{Le jeu du Mathador}

Présentation issue du site \url{https://www.mathador.fr}:

\begin{tcolorbox}[enhanced,colback=orange!15,colframe=orange!75!black,fontupper=\small\itshape\sffamily,width=0.75\linewidth,center,watermark opacity=0.1,watermark zoom=0.95,watermark graphics={mathadorlogo.pdf}]
\textbf{Le principe : lancez les dés !}

Ils vous donnent 1 nombre-cible (entre 0 et 99) et 5 nombres à utiliser pour l'atteindre.

Vous avez les 4 opérations pour y arriver, mais attention : le temps est limité et les opérations rapportent des points\ldots

Ces 6 dés donnent des milliers de \textit{compte-est-bon}. Pour chacun, jusqu'à des centaines de réponses différentes : de quoi se remuer les méninges et chercher la stratégie qui vous fera gagner !
\end{tcolorbox}

Les éléments graphiques ont été repris du site officiel, et adaptés/transcrits en langage \TikZ.

\subsection{Crédits, mentions légales}

\og \textsf{Mathador} est une marque protégée de Monsieur Eric Trouillot et de Réseau Canopé.

Eric Trouillot est le concepteur du jeu \textsf{Mathador} dont vous pouvez retrouver toute l’actualité sur le site \url{www.mathador.fr}.

Ce package \LaTeX\ est une ressource indépendante non affiliée à \url{www.mathador.fr}. \fg

\subsection{Chargement, packages utiles}

Pour charger le package \texttt{mathador}, il suffit de placer, dans le préambule :

\begin{DemoCode}{listing only}
%version 'normale'
\usepackage{mathador}
%version avec le package bm pour les opérateurs en gras
\usepackage[bm]{mathador}
%version sans la librairie tikz.babel
\usepackage[nonbabel]{mathador}
\end{DemoCode}

Les packages utilisés sont \ShowCode{tikz}, \ShowCode{ifthen}, \ShowCode{listofitems}, \ShowCode{simplekv} et éventuellement \ShowCode{bm}.

\subsection{Idées globales}

Les éléments textuels sont automatiquement adaptés aux dimensions du plateau et des cases.

Les styles sont fixés, mais peuvent être redéfinis

\begin{DemoCode}{listing only}
%====COULEURS PRINCIPALES
\definecolor{mathadorppal}{HTML}{392b58}
\definecolor{mathadorsecond}{HTML}{4b3971}
\definecolor{mathadorjaune}{HTML}{dfc447}
\definecolor{mathadorcible}{HTML}{e1c444}
\definecolor{mathadorrose}{HTML}{c25a75}
\definecolor{mathadorvert}{HTML}{aaa640}
\definecolor{mathadorbleu}{HTML}{69acbd}
%====COULEURS PASTELS
\definecolor{mathadorppalpastel}{HTML}{D9E8FF}
\definecolor{mathadorsecondpastel}{HTML}{EDF4FF}
%====STYLES TIKZ
\tikzset{mathadorpolice/.style={font=\sffamily}}
\tikzset{mathadorpolicegras/.style={font=\bfseries\sffamily}}
\end{DemoCode}

Les opérateurs mathématiques sont, par défaut, en \ShowCode{\\mathsf}, mais il peuvent être redéfinis.

\begin{DemoCode}{listing only}
%====OPÉRATEURS
\newcommand\intmathadorgrassf[1]{%
	\@ifpackageloaded{bm}{$\bm{\mathsf{#1}}$}{$\mathsf{#1}$}%
}
\def\mathadorplussymb{\intmathadorgrassf{\vphantom{-\times\div+}+}}
\def\mathadormoinssymb{\intmathadorgrassf{\vphantom{-\times\div+}-}}
\def\mathadorfoissymb{\intmathadorgrassf{\vphantom{-\times\div+}\times}}
\def\mathadordivsymb{\intmathadorgrassf{\vphantom{-\times\div+}\div}}
\end{DemoCode}

\section{Les commandes}

\subsection{Le plateau}

La commande pour créer le plateau de jeu est :

\begin{DemoCode}{listing only}
\PlateauMathador[clés]<options tikz>{valeurs (5)}{valeur cible}
\end{DemoCode}

Les \textsf{clés} disponibles sont :

\begin{itemize}
	\item \ShowCode{rayon} : rayon du plateau\hfill(\ShowCode{5cm} par défaut)
	\item \ShowCode{logo} : booléen pour afficher le logo\hfill(\ShowCode{true} par défaut)
	\item \ShowCode{epaisseur} : épaisseur (\texttt{auto} ou \texttt{longueur})\hfill(\ShowCode{auto} par défaut)
	\item \ShowCode{ombre} : booléen pour afficher une ombre\hfill(\ShowCode{true} par défaut)
	\item \ShowCode{impression} : booléen pour un mode NB\hfill(\ShowCode{false} par défaut)
	\item \ShowCode{pastels} : booléen pour un mode pastel \textit{bleu} si \ShowCode{impression=true} \hfill(\ShowCode{false} par défaut)
\end{itemize}

\begin{DemoCode}{}
\PlateauMathador{1,6,19,7,12}{39}
\end{DemoCode}

\begin{DemoCode}{}
\PlateauMathador[rayon=2.5cm,logo=false,epaisseur=0.2mm]{1,6,19,7,12}{39}
\end{DemoCode}

\begin{DemoCode}{}
\PlateauMathador[impression,rayon=2.5cm]{1,6,19,7,12}{39}
\end{DemoCode}

\begin{DemoCode}{}
\PlateauMathador[impression,pastels,rayon=2.5cm]{1,6,19,7,12}{39}
\end{DemoCode}

\begin{DemoCode}{}
%opérateurs basés sur fontawesome
\def\mathadorplussymb{\vphantom{\faPlus\faMinus\faDivide\faTimes}\faPlus}
\def\mathadormoinssymb{\vphantom{\faPlus\faMinus\faDivide\faTimes}\faMinus}
\def\mathadorfoissymb{\vphantom{\faPlus\faMinus\faDivide\faTimes}\faTimes}
\def\mathadordivsymb{\vphantom{\faPlus\faMinus\faDivide\faTimes}\faDivide}

\PlateauMathador[rayon=1.75cm,logo=false]{,,,,}{}
\end{DemoCode}

\subsection{La cible, les dés}

Il est également possible de présenter les valeurs via les 6 \textit{dés} :

\begin{DemoCode}{listing only}
%la cible
\CibleMathador[couleur]<options tikz>{hauteur}{valeur}
\end{DemoCode}

\begin{DemoCode}{listing only}
%les dés
\DeMathador<options tikz>{couleur}{type}{hauteur}{numéro}
\end{DemoCode}

La hauteur des dés n'est pas forcément uniforme, l'idée est de pouvoir présenter les dés \textit{cote à cote}, avec un rendu satisfaisant.

\begin{DemoCode}{}
\DeMathador{red}{Q}{2cm}{4}             %dé Quatre
\DeMathador{olive!50!lime}{S}{2cm}{4}   %dé Six
\DeMathador{orange}{H}{2cm}{4}          %dé Huit
\DeMathador{blue}{D}{2cm}{4}            %dé Douze
\DeMathador{purple}{V}{2cm}{4}          %dé Vingt
\end{DemoCode}

\begin{DemoCode}{}
\CibleMathador{2cm}{61}
\end{DemoCode}

\begin{DemoCode}{}
\CibleMathador[teal!50!black]{1.5cm}{61}
\end{DemoCode}

\subsection{Tirage complet}

L'idée est de proposer une commande pour présenter un tirage \textit{complet}, avec les 7 valeurs.

\begin{DemoCode}{listing only}
\TirageDesMathador[clés]{cible}{dé1}{dé2}{dé3}{dé4}{dé5}
\end{DemoCode}

Les \textsf{clés} disponibles sont :

\begin{itemize}
	\item \ShowCode{hauteur} : hauteur des dés \hfill(\ShowCode{1.5cm} par défaut)
	\item \ShowCode{couleurs} : couleurs des dés\hfill(\ShowCode{purple!50!black,red,olive!50!lime,orange,blue,purple} par défaut)
	\item \ShowCode{superpos} : booléen pour \textit{superposer} les dés\hfill(\ShowCode{true} par défaut)
	\item \ShowCode{impression} : booléen pour un mode NB\hfill(\ShowCode{false} par défaut)
\end{itemize}

\begin{DemoCode}{}
\TirageDesMathador{2}{4}{6}{7}{9}{61}
\end{DemoCode}

\begin{DemoCode}{}
\TirageDesMathador[impression]{2}{4}{6}{7}{9}{61}
\end{DemoCode}

\pagebreak

\section{Historique}

\texttt{0.1.0 : Meilleure gestion des opérateurs}

\texttt{0.1.0 : Version initiale}
%
%\section{The code}
%
%\DemoCodeFile{commalists-tools-l3.sty}

\end{document}
